%%%%%%%%%%%%%%%%%%%%%%%%%%%%%%%%%%%%%%%%%%%%%%%%%%%%%%%%%%%%%%%%%%%%%%%%%%%%%
%% PREAMBLE
%% Reorganized for readability: packages, colors, custom commands, code
%% listing styles, tcolorbox theorem-like environments and title formatting
%% are grouped by purpose. No functional change was made to the original
%% document: only the order of declarations and the comments were edited.
%%%%%%%%%%%%%%%%%%%%%%%%%%%%%%%%%%%%%%%%%%%%%%%%%%%%%%%%%%%%%%%%%%%%%%%%%%%%%


%%%%%%%%%%%%%%%%%%%%%%%%%%%%%%%%%%%%%%%%%%%%%%%%%%%%%%%%%%%%%%%%%%%%%%%%%%%%%
%% 1. PACKAGES
%%%%%%%%%%%%%%%%%%%%%%%%%%%%%%%%%%%%%%%%%%%%%%%%%%%%%%%%%%%%%%%%%%%%%%%%%%%%%

% --- Encoding, language and fonts ---------------------------------------
\usepackage[utf8]{inputenc}
\usepackage[T1]{fontenc}
\usepackage[english]{babel}
\usepackage{mathpazo}          % Palatino font, replaces lmodern to match the image's serif style
\usepackage[scaled]{helvet}    % Helvetica, used for section/subsection titles

% --- Page layout and headers/footers ------------------------------------
\usepackage[paper=a4paper,left=25mm,right=25mm,bottom=25mm,top=25mm]{geometry}
\usepackage{fancyhdr}
\usepackage[Rejne]{fncychap}
\usepackage{titlesec}
\usepackage{titletoc}
\usepackage{needspace}

% --- Math ----------------------------------------------------------------
\usepackage{amsmath}
\usepackage{amssymb}
\usepackage{amsfonts}
\usepackage{mathtools}
\usepackage{mathrsfs}
\usepackage{stmaryrd}
\usepackage{nicefrac, xfrac}

% --- Graphics and TikZ -----------------------------------------------------
\usepackage{graphicx}          % Required for inserting images
\usepackage[export]{adjustbox}
\usepackage{pdfpages}
\usepackage{tikz}
\usepackage{tikz-3dplot}
\usepackage{tikz-cd}
\usepackage{pgfplots}
\usepackage{pgf-pie}
\usetikzlibrary{positioning}

% --- Colors and colored boxes ---------------------------------------------
\usepackage[table,xcdraw]{xcolor}
\usepackage[most]{tcolorbox}
\tcbuselibrary{theorems}

% --- Lists, layout helpers and decorations --------------------------------
\usepackage[shortlabels]{enumitem}
\usepackage{wrapfig}
\usepackage{multicol}
\usepackage{moresize}
\usepackage{fourier-orns}       % Decorative ornaments (used by \flowerLine, \eqImportante)
\usepackage{psvectorian}        % Decorative vectorial ornaments
\usepackage{lipsum}             % Placeholder ("dummy") text
\usepackage{blindtext}          % Placeholder ("dummy") text

% --- Code listings ---------------------------------------------------------
\usepackage{listings}
\usepackage{minted}

% --- Miscellaneous ---------------------------------------------------------
\usepackage{hyperref}


%%%%%%%%%%%%%%%%%%%%%%%%%%%%%%%%%%%%%%%%%%%%%%%%%%%%%%%%%%%%%%%%%%%%%%%%%%%%%
%% 2. COLOR DEFINITIONS
%%%%%%%%%%%%%%%%%%%%%%%%%%%%%%%%%%%%%%%%%%%%%%%%%%%%%%%%%%%%%%%%%%%%%%%%%%%%%

% --- General purpose colors ---
\definecolor{light-gray}{gray}{0.95}
\definecolor{darkgreen}{HTML}{008000}
\definecolor{g}{RGB}{60, 50, 50}

% --- "Sapienza" theme colors ---
\definecolor{sapienza}{HTML}{660f1d}
\definecolor{lightSapienza}{HTML}{e3d3d5}

% --- Cover / decorative colors ---
\definecolor{cop}{HTML}{f7ecd7}
\definecolor{copAut}{HTML}{ababab}
\definecolor{copAut2}{HTML}{c3c3e6}
\definecolor{purcop}{HTML}{d0d3db}
\definecolor{cartaRiciclata}{HTML}{fcfcf7}

% --- Colors used for C / C++ code listings ---
\definecolor{mGreen}{rgb}{0,0.6,0}
\definecolor{mGray}{rgb}{0.5,0.5,0.5}
\definecolor{mPurple}{rgb}{0.58,0,0.82}
\definecolor{backgroundColour}{rgb}{0.95,0.95,0.92}


%%%%%%%%%%%%%%%%%%%%%%%%%%%%%%%%%%%%%%%%%%%%%%%%%%%%%%%%%%%%%%%%%%%%%%%%%%%%%
%% 3. CUSTOM COMMANDS
%%%%%%%%%%%%%%%%%%%%%%%%%%%%%%%%%%%%%%%%%%%%%%%%%%%%%%%%%%%%%%%%%%%%%%%%%%%%%

% --- Colored text shortcuts ------------------------------------------------
\newcommand{\redText}[1]{\color{red}#1\color{black}}
\newcommand{\blueText}[1]{\color{blue}#1\color{black}}
\newcommand{\greenText}[1]{\color{darkgreen}#1\color{black}}
\newcommand{\sapText}[1]{\color{sapienza}#1\color{black}}
\newcommand{\textg}[1]{\color{g}{\textbf{#1}}\color{black}}

% --- Inline code / shell snippets -------------------------------------------
\newcommand{\code}[1]{\colorbox{light-gray}{\texttt{#1}}}
\newcommand{\codee}[1]{\colorbox{white}{\texttt{#1}}}
\newcommand{\shell}[1]{\colorbox{black}{\textcolor{white}{\texttt{#1}}}}
\newcommand{\boxedMath}[1]{\begin{tabular}{|c|}\hline \texttt{#1} \\ \hline\end{tabular} :}

% --- Highlighted paragraph boxes --------------------------------------------
\newcommand\greybox[1]{%
  \vskip\baselineskip%
  \par\noindent\colorbox{light-gray}{%
    \begin{minipage}{\textwidth}#1\end{minipage}%
  }%
  \vskip\baselineskip%
}
\newcommand\sapbox[1]{%
  \vskip\baselineskip%
  \par\noindent\colorbox{lightSapienza}{%
    \begin{minipage}{\textwidth}#1\end{minipage}%
  }%
  \vskip\baselineskip%
}

% --- Statement labels (theorem/definition/etc. headers as plain text) ------
\newcommand{\teo}[1]{{\large\color{sapienza}\textbf{Teorema #1 :\hphantom{a}}}}
\newcommand{\defi}[1]{{\large\color{sapienza}\textbf{Definizione #1 :\hphantom{a}}}}
\newcommand{\prop}[1]{{\color{sapienza}\textbf{Proposizione #1 :\hphantom{a}}}}
\newcommand{\lemma}[1]{{\color{sapienza}\textbf{Lemma #1 :\hphantom{a}}}}
\newcommand{\claim}[1]{{\color{sapienza}\textbf{Claim #1 :\hphantom{a}}}}
\newcommand{\dimo}[1]{{\color{sapienza}\textbf{Dimostrazione #1 :\hphantom{a}}}}
\newcommand{\acc}{\\\hphantom{}\\}
\newcommand{\esempio}[1]{{\acc\large\color{sapienza}\textbf{Esempio #1 \hphantom{a}}\acc}}

% --- Number sets and blackboard-bold symbols --------------------------------
\newcommand{\N}{{\mathbb N}}
\newcommand{\Z}{{\mathbb Z}}
\newcommand{\R}{{\mathbb R}}
\newcommand{\C}{{\mathbb C}}
\newcommand{\K}{{\mathbb K}}
\newcommand{\quat}{{\mathbb H}}
\newcommand{\Prob}{{\mathbb P}}
\newcommand{\VA}{{\mathbb E}}

% --- Calligraphic / script symbols ------------------------------------------
\newcommand{\Sn}{{\mathcal S_n}}
\newcommand{\An}{{\mathcal A_n}}
\newcommand{\E}{{\mathcal E}}
\newcommand{\B}{{\mathcal B}}

% --- Bold vector/greek symbols ----------------------------------------------
\newcommand{\btheta}{{\boldsymbol{\theta}}}
\newcommand{\btau}{{\boldsymbol{\tau}}}
\newcommand{\bomega}{{\boldsymbol{\omega}}}
\newcommand{\Bomega}{{\boldsymbol{\Omega}}}
\newcommand{\bmu}{{\boldsymbol{\mu}}}
\newcommand{\brho}{{\boldsymbol{\rho}}}
\newcommand{\bnu}{{\boldsymbol{\nu}}}
\newcommand{\bsigma}{{\boldsymbol{\sigma}}}
\newcommand{\blambda}{{\boldsymbol{\lambda}}}

% --- Dotted mechanics symbols (velocity/acceleration notation) -------------
\newcommand{\dq}{{\dot{\mathbf q}}}
\newcommand{\de}{{\dot{\mathbf e}}}
\newcommand{\dy}{{\dot{\mathbf y}}}
\newcommand{\ddq}{{\ddot{\mathbf q}}}
\newcommand{\ddy}{{\ddot{\mathbf y}}}
\newcommand{\ve}{{\bar v}}

% --- Group theory / algebra operators ---------------------------------------
\newcommand{\aut}{{\text{Aut}}}
\newcommand{\Span}{{\text{Span}}}
\newcommand{\End}{{\text{End}}}
\newcommand{\cen}{{\text{Centro}}}
\newcommand{\norm}{{\unlhd}}
\newcommand{\ciclS}{{\left \langle }}
\newcommand{\ciclE}{{\right \rangle }}
\newcommand{\bra}[1]{\langle #1 \rangle}
\newcommand{\mcm}{{\text{mcm}}}
\newcommand{\MCD}{{\text{MCD}}}
\newcommand{\rg}{{\text{rg}}}
\newcommand{\supp}{{\text{Supp}}}

% --- Complexity theory symbols -----------------------------------------------
\newcommand{\ridFunc}{{f:\Sigma^*\rightarrow \Sigma^*}}
\newcommand{\rid}{{\le_m^P}}
\newcommand{\blank}{{\sqcup}}

% --- Miscellaneous math symbols and utilities -------------------------------
\newcommand{\quadBrack}[1]{ \llbracket #1 \rrbracket }
\newcommand{\notimplies}{%
  \mathrel{{\ooalign{\hidewidth$\not\phantom{=}$\hidewidth\cr$\implies$}}}}
\newcommand{\tc}{{\text{ tale che }}}
\newcommand{\spaz}{{\text{\hphantom{aa}}}}

% --- Decorative separators ---------------------------------------------------
\newcommand{\flowerLine}{ \begin{center}\decofourleft\hphantom{ }\decoone\hphantom{ }\decofourright\hphantom{}\hphantom{aa}
\decofourleft\hphantom{ }\decoone\hphantom{ }\decofourright\hphantom{}\hphantom{aa}
\decofourleft\hphantom{ }\decoone\hphantom{ }\decofourright\hphantom{}\hphantom{aa}
\decofourleft\hphantom{ }\decoone\hphantom{ }\decofourright\hphantom{}\hphantom{aa} 
\decofourleft\hphantom{ }\decoone\hphantom{ }\decofourright\hphantom{}\hphantom{aa}
\decofourleft\hphantom{ }\decoone\hphantom{ }\decofourright\hphantom{}\hphantom{aa}
\decofourleft\hphantom{ }\decoone\hphantom{ }\decofourright\hphantom{}\hphantom{aa}
\decofourleft\hphantom{ }\decoone\hphantom{ }\decofourright\hphantom{}\hphantom{aa}
\decofourleft\hphantom{ }\decoone\hphantom{ }\decofourright\hphantom{}\hphantom{aa}
\end{center}}
\newcommand{\eqImportante}[1]{\begin{center}\huge\lefthand\hphantom{a}
    \normalsize\texttt{#1}
    \hphantom{aaa}\huge\righthand\end{center}}


%%%%%%%%%%%%%%%%%%%%%%%%%%%%%%%%%%%%%%%%%%%%%%%%%%%%%%%%%%%%%%%%%%%%%%%%%%%%%
%% 4. PAGE STYLE (HEADER / FOOTER)
%%%%%%%%%%%%%%%%%%%%%%%%%%%%%%%%%%%%%%%%%%%%%%%%%%%%%%%%%%%%%%%%%%%%%%%%%%%%%
\fancyhf{}
\pagestyle{fancy}
\renewcommand{\headrulewidth}{0.4pt}


%%%%%%%%%%%%%%%%%%%%%%%%%%%%%%%%%%%%%%%%%%%%%%%%%%%%%%%%%%%%%%%%%%%%%%%%%%%%%
%% 5. CODE LISTING STYLES (C / C++)
%%%%%%%%%%%%%%%%%%%%%%%%%%%%%%%%%%%%%%%%%%%%%%%%%%%%%%%%%%%%%%%%%%%%%%%%%%%%%
% The following styles configure the `listings` package for displaying
% C and C++ source code with syntax highlighting.

\lstdefinestyle{CStyle}{
    backgroundcolor=\color{backgroundColour},   
    commentstyle=\color{mGreen},
    keywordstyle=\color{magenta},
    numberstyle=\tiny\color{mGray},
    stringstyle=\color{mPurple},
    basicstyle=\footnotesize,
    breakatwhitespace=false,         
    breaklines=true,                 
    captionpos=b,                    
    keepspaces=true,                 
    numbers=left,                    
    numbersep=5pt,                  
    showspaces=false,                
    showstringspaces=false,
    showtabs=false,                  
    tabsize=2,
    language=C
}

\lstdefinestyle{CppStyle}{
    backgroundcolor=\color{backgroundColour},   
    commentstyle=\color{mGreen}\ttfamily,
    morecomment=[l][\color{magenta}]{\#}
    keywordstyle=\color{blue}\ttfamily,
    numberstyle=\tiny\color{mGray},
    stringstyle=\color{red}\ttfamily,
    basicstyle=\ttfamily,
    breakatwhitespace=false,         
    breaklines=true,                 
    captionpos=b,                    
    keepspaces=true,                 
    numbers=left,                    
    numbersep=5pt,                  
    showspaces=false,                
    showstringspaces=false,
    showtabs=false,                  
    tabsize=2,
    language=C
}

% Global default `listings` configuration (applied when no style is selected)
\lstset{language=C++,
                basicstyle=\ttfamily,
                keywordstyle=\color{blue}\ttfamily,
                stringstyle=\color{red}\ttfamily,
                commentstyle=\color{green}\ttfamily,
                morecomment=[l][\color{magenta}]{\#}
}


%%%%%%%%%%%%%%%%%%%%%%%%%%%%%%%%%%%%%%%%%%%%%%%%%%%%%%%%%%%%%%%%%%%%%%%%%%%%%
%% 6. TCOLORBOX THEOREM-LIKE ENVIRONMENTS
%%%%%%%%%%%%%%%%%%%%%%%%%%%%%%%%%%%%%%%%%%%%%%%%%%%%%%%%%%%%%%%%%%%%%%%%%%%%%
% Each environment below renders a boxed, numbered statement (definition,
% theorem, proposition, ...) with a "tab" style title in the top-left
% corner, following the same visual template.

\newtcbtheorem[number within=section]{boxdefinition}{Definition}{
    enhanced,
    before skip=15pt, after skip=15pt,
    colback=white,                          % white background for the text
    colframe=sapienza!50!black!30,          % border color (light green/gray)
    fonttitle=\bfseries,
    coltitle=black,
    % Style of the top-left "tab"
    attach boxed title to top left={xshift=10mm, yshift=-3.5mm, yshifttext=-1mm},
    boxed title style={
        colframe=sapienza!50!black!30,
        arc=3mm,
        outer arc=3mm,
        left=10pt, right=10pt,
        top=2pt, bottom=2pt,
        % Green gradient shading, as in the reference image
        interior style={left color=sapienza!20, right color=sapienza!40}
    },
    % Rounding of the main box
    arc=4mm,
    separator sign none,                    % removes the automatic colon so it can be added manually in the title
    description font=\normalfont\bfseries,
    description delimiters={: }{ }          % formats the title as "Definition 2.1: Title"
}{def}

\newtcbtheorem[number within=section]{boxtheorem}{Theorem}{
    enhanced,
    before skip=15pt, after skip=15pt,
    colback=white,
    colframe=sapienza!50!black!30,
    fonttitle=\bfseries,
    coltitle=black,
    attach boxed title to top left={xshift=10mm, yshift=-3.5mm, yshifttext=-1mm},
    boxed title style={
        colframe=sapienza!50!black!30,
        arc=3mm,
        outer arc=3mm,
        left=10pt, right=10pt,
        top=2pt, bottom=2pt,
        interior style={left color=sapienza!20, right color=sapienza!40}
    },
    arc=4mm,
    separator sign none,
    description font=\normalfont\bfseries,
    description delimiters={: }{ }
}{def}

\newtcbtheorem[number within=section]{boxproposition}{Proposition}{
    enhanced,
    before skip=15pt, after skip=15pt,
    colback=white,
    colframe=sapienza!50!black!30,
    fonttitle=\bfseries,
    coltitle=black,
    attach boxed title to top left={xshift=10mm, yshift=-3.5mm, yshifttext=-1mm},
    boxed title style={
        colframe=sapienza!50!black!30,
        arc=3mm,
        outer arc=3mm,
        left=10pt, right=10pt,
        top=2pt, bottom=2pt,
        interior style={left color=sapienza!30, right color=sapienza!50}
    },
    arc=4mm,
    separator sign none,
    description font=\normalfont\bfseries,
    description delimiters={: }{ }
}{def}

\newtcbtheorem[number within=section]{boxequation}{Equation}{
    enhanced,
    before skip=15pt, after skip=15pt,
    colback=white,
    colframe=sapienza!50!black!30,
    fonttitle=\bfseries,
    coltitle=black,
    attach boxed title to top left={xshift=10mm, yshift=-3.5mm, yshifttext=-1mm},
    boxed title style={
        colframe=sapienza!50!black!30,
        arc=3mm,
        outer arc=3mm,
        left=10pt, right=10pt,
        top=2pt, bottom=2pt,
        interior style={left color=sapienza!30, right color=sapienza!50}
    },
    arc=4mm,
    separator sign none,
    description font=\normalfont\bfseries,
    description delimiters={: }{ }
}{def}

\newtcbtheorem[number within=section]{boxlemma}{Lemma}{
    enhanced,
    before skip=15pt, after skip=15pt,
    colback=white,
    colframe=sapienza!50!black!30,
    fonttitle=\bfseries,
    coltitle=black,
    attach boxed title to top left={xshift=10mm, yshift=-3.5mm, yshifttext=-1mm},
    boxed title style={
        colframe=sapienza!50!black!30,
        arc=3mm,
        outer arc=3mm,
        left=10pt, right=10pt,
        top=2pt, bottom=2pt,
        interior style={left color=sapienza!30, right color=sapienza!50}
    },
    arc=4mm,
    separator sign none,
    description font=\normalfont\bfseries,
    description delimiters={: }{ }
}{def}

\newtcbtheorem[number within=section]{boxobservation}{Observation}{
    enhanced,
    before skip=15pt, after skip=15pt,
    colback=white,
    colframe=sapienza!50!black!30,
    fonttitle=\bfseries,
    coltitle=black,
    attach boxed title to top left={xshift=10mm, yshift=-3.5mm, yshifttext=-1mm},
    boxed title style={
        colframe=sapienza!50!black!30,
        arc=3mm,
        outer arc=3mm,
        left=10pt, right=10pt,
        top=2pt, bottom=2pt,
        interior style={left color=sapienza!30, right color=sapienza!50}
    },
    arc=4mm,
    separator sign none,
    description font=\normalfont\bfseries,
    description delimiters={: }{ }
}{def}

\newtcbtheorem[number within=section]{boxpostulate}{Postulate}{
    enhanced,
    before skip=15pt, after skip=15pt,
    colback=white,
    colframe=sapienza!50!black!30,
    fonttitle=\bfseries,
    coltitle=black,
    attach boxed title to top left={xshift=10mm, yshift=-3.5mm, yshifttext=-1mm},
    boxed title style={
        colframe=sapienza!50!black!30,
        arc=3mm,
        outer arc=3mm,
        left=10pt, right=10pt,
        top=2pt, bottom=2pt,
        interior style={left color=sapienza!30, right color=sapienza!50}
    },
    arc=4mm,
    separator sign none,
    description font=\normalfont\bfseries,
    description delimiters={: }{ }
}{def}


%%%%%%%%%%%%%%%%%%%%%%%%%%%%%%%%%%%%%%%%%%%%%%%%%%%%%%%%%%%%%%%%%%%%%%%%%%%%%
%% 7. SECTION / SUBSECTION TITLE FORMATTING
%%%%%%%%%%%%%%%%%%%%%%%%%%%%%%%%%%%%%%%%%%%%%%%%%%%%%%%%%%%%%%%%%%%%%%%%%%%%%

% Section titles: Helvetica, large bold number, colored rule underneath
\titleformat{\section}
  {\needspace{3\baselineskip} % requires at least 3 lines of free space before the title
   \fontfamily{phv}\selectfont\huge\bfseries} 
  {\color{sapienza}\fontsize{20}{20}\selectfont\thesection} 
  {12pt} 
  {} 
  [\vspace{2pt}\nopagebreak\color{sapienza!40}\hrule height 3pt]

% Subsection titles: same style, slightly smaller
\titleformat{\subsection}
  {\needspace{2\baselineskip} % requires at least 2 lines of free space before the title
   \fontfamily{phv}\selectfont\Large\bfseries} 
  {\color{sapienza}\fontsize{16}{16}\selectfont\thesubsection} 
  {12pt} 
  {} 
  [\vspace{2pt}\nopagebreak\color{sapienza!40}\hrule height 1.5pt]